% problem_id: S001-lemma-trivial-fibres-dvr
% source_id: S001 (Stacks Project)
% source_file: flat.tex
% statement_locator: flat.tex lines 12485--12496
% proof_locator: flat.tex lines 12498--12638
% env: lemma
% id_note: label 在本来源内唯一
% pairing: tight (verified)
% status: complete (statement + author proof verbatim + reason + locators)
%
\begin{lemma}
\label{lemma-trivial-fibres-dvr}
Let $f : X \to S$ be a proper morphism with geometrically connected fibres
where $S$ is the spectrum of a discrete valuation ring. Denote $\eta \in S$
the generic point and denote $X_n \subset X$ the closed subscheme
cutout by the $n$th power of a uniformizer on $S$.
Then there exists
an integer $n$ such that the following is true: any finite
locally free $\mathcal{O}_X$-module $\mathcal{E}$
such that $\mathcal{E}|_{X_\eta}$ and $\mathcal{E}|_{X_n}$
are free, is free.
\end{lemma}

\begin{proof}
We first reduce to the case where $X \to S$ has a section. Say $S = \Spec(A)$.
Choose a closed point $\xi$ of $X_\eta$. Choose an extension
of discrete valuation rings $A \subset B$ such that the fraction field
of $B$ is $\kappa(\xi)$. This is possible by Krull-Akizuki
(Algebra, Lemma \ref{algebra-lemma-integral-closure-Dedekind})
and the fact that $\kappa(\xi)$ is a finite extension of the
fraction field of $A$.
By the valuative criterion of properness
(Morphisms, Lemma \ref{morphisms-lemma-characterize-proper})
we get a $B$-valued point $\tau : \Spec(B) \to X$
which induces a section $\sigma : \Spec(B) \to X_B$.
For a finite locally free $\mathcal{O}_X$-module $\mathcal{E}$
let $\mathcal{E}_B$ be the pullback to the base change $X_B$.
By flat base change
(Cohomology of Schemes, Lemma \ref{coherent-lemma-flat-base-change-cohomology})
we see that $H^0(X_B, \mathcal{E}_B) = H^0(X, \mathcal{E}) \otimes_A B$.
Thus if $\mathcal{E}_B$ is free of rank $r$, then the sections in
$H^0(X, \mathcal{E})$ generate the free $B$-module
$\tau^*\mathcal{E} = \sigma^*\mathcal{E}_B$.
In particular, we can find $r$ global sections $s_1, \ldots, s_r$
of $\mathcal{E}$ which generate $\tau^*\mathcal{E}$. Then
$$
s_1, \ldots, s_r :
\mathcal{O}_X^{\oplus r}
\longrightarrow
\mathcal{E}
$$
is a map of finite locally free $\mathcal{O}_X$-modules of rank $r$
and the pullback to $X_B$ is a map of free $\mathcal{O}_{X_B}$-modules
which restricts to an isomorphism in one point of each fibre.
Taking the determinant we get a function
$g \in \Gamma(X_\eta, \mathcal{O}_{X_B})$
which is invertible in one point of each fibre.
As the fibres are proper and connected, we see that $g$
must be invertible (details omitted; hint: use Varieties, Lemma
\ref{varieties-lemma-proper-geometrically-reduced-global-sections}).
Thus it suffices to prove the lemma for the base change $X_B \to \Spec(B)$.

\medskip\noindent
Assume we have a section $\sigma : S \to X$. Let $\mathcal{E}$
be a finite locally free $\mathcal{O}_X$-module which is assumed
free on the generic fibre and on $X_n$ (we will choose $n$ later).
Choose an isomorphism $\sigma^*\mathcal{E} = \mathcal{O}_S^{\oplus r}$.
Consider the map
$$
K = R\Gamma(X, \mathcal{E}) \longrightarrow
R\Gamma(S, \sigma^*\mathcal{E}) = A^{\oplus r}
$$
in $D(A)$. Arguing as above, we see $\mathcal{E}$ is free if (and only if)
the induced map $H^0(K) = H^0(X, \mathcal{E}) \to A^{\oplus r}$ is surjective.

\medskip\noindent
Set $L = R\Gamma(X, \mathcal{O}_X^{\oplus r})$ and observe that the
corresponding map $L \to A^{\oplus r}$ has the desired property.
Observe that $K \otimes_A Q(A) \cong L \otimes_A Q(A)$
by flat base change and the assumption that $\mathcal{E}$
is free on the generic fibre. Let $\pi \in A$ be a uniformizer. Observe that
$$
K \otimes_A^\mathbf{L} A/\pi^m A =
R\Gamma(X, \mathcal{E} \xrightarrow{\pi^m} \mathcal{E})
$$
and similarly for $L$.
Denote $\mathcal{E}_{tors} \subset \mathcal{E}$ the coherent subsheaf of
sections supported on the special fibre and similarly for other
$\mathcal{O}_X$-modules. Choose $k > 0$ such that
$(\mathcal{O}_X)_{tors} \to \mathcal{O}_X/\pi^k \mathcal{O}_X$
is injective (Cohomology of Schemes, Lemma \ref{coherent-lemma-Artin-Rees}).
Since $\mathcal{E}$ is locally free, we see
that $\mathcal{E}_{tors} \subset \mathcal{E}/\pi^k\mathcal{E}$.
Then for $n \geq m + k$ we have isomorphisms
\begin{align*}
(\mathcal{E} \xrightarrow{\pi^m} \mathcal{E})
& \cong
(\mathcal{E}/\pi^k\mathcal{E} \xrightarrow{\pi^m}
\mathcal{E}/\pi^{k + m}\mathcal{E}) \\
& \cong
(\mathcal{O}_X^{\oplus r}/\pi^k\mathcal{O}_X^{\oplus r} \xrightarrow{\pi^m}
\mathcal{O}_X^{\oplus r}/\pi^{k + m}\mathcal{O}_X^{\oplus r}) \\
& \cong
(\mathcal{O}_X^{\oplus r} \xrightarrow{\pi^m} \mathcal{O}_X^{\oplus r})
\end{align*}
in $D(\mathcal{O}_X)$. This determines an isomorphism
$$
K \otimes_A^\mathbf{L} A/\pi^m A \cong L \otimes_A^\mathbf{L} A/\pi^m A
$$
in $D(A)$ (holds when $n \geq m + k$). Observe that these isomorphisms
are compatible with pulling back by $\sigma$ hence in particular
we conclude that
$K \otimes_A^\mathbf{L} A/\pi^m A \to (A/\pi^m A)^{\oplus r}$
defines an surjection on degree $0$ cohomology modules (as
this is true for $L$).
Since $A$ is a discrete valuation ring, we have
$$
K \cong \bigoplus H^i(K)[-i]
\quad\text{and}\quad
L \cong
\bigoplus H^i(L)[-i]
$$
in $D(A)$. See More on Algebra, Example
\ref{more-algebra-example-finite-injective-finite-global-dimension}.
The cohomology groups $H^i(K) = H^i(X, \mathcal{E})$ and
$H^i(L) = H^i(X, \mathcal{O}_X)^{\oplus r}$
are finite $A$-modules by Cohomology of Schemes, Lemma
\ref{coherent-lemma-proper-over-affine-cohomology-finite}.
By More on Algebra, Lemma
\ref{more-algebra-lemma-generalized-valuation-ring-modules}
these modules are direct sums of cyclic modules.
We have seen above that the rank $\beta_i$ of the
free part of $H^i(K)$ and $H^i(L)$ are the same.
Next, observe that
$$
H^i(L \otimes_A^\mathbf{L} A/\pi^m A) =
H^i(L)/\pi^m H^i(L) \oplus H^{i + 1}(L)[\pi^m]
$$
and similarly for $K$. Let $e$ be the largest integer
such that $A/\pi^eA$ occurs as a summand of $H^i(X, \mathcal{O}_X)$,
or equivalently $H^i(L)$, for some $i$. Then taking $m = e + 1$
we see that $H^i(L \otimes_A^\mathbf{L} A/\pi^m A)$ is a direct sum of
$\beta_i$ copies of $A/\pi^m A$ and some other cyclic modules
each annihilated by $\pi^e$. By the same reasoning for $K$
and the isomorphism
$K \otimes_A^\mathbf{L} A/\pi^m A \cong L \otimes_A^\mathbf{L} A/\pi^m A$
it follows that $H^i(K)$
cannot have any cyclic summands of the form $A/\pi^l A$
with $l > e$. (It also follows that $K$ is isomorphic to $L$
as an object of $D(A)$, but we won't need this.)
Then the only way the map
$$
H^0(K \otimes^\mathbf{L}_A A/\pi^{e + 1} A) =
H^0(K)/\pi^{e + 1}H^0(K) \oplus H^1(K)[\pi^{e + 1}]
\longrightarrow
(A/\pi^{e + 1} A)^{\oplus r}
$$
is surjective, is if it is surjective on the
first summand. This is what we wanted to show.
(To be precise, the integer $n$ in the statement of
the lemma, if there is a section $\sigma$,
should be equal to $k + e + 1$ where $k$ and $e$ are as above
and depend only on $X$.)
\end{proof}
